% Section paragraphs written during the image check; citations identify typescript pages.
\subsection*{14. Homotopical algebra and the limits of present topology (pp. 1--2)}
Grothendieck opens the continuation of the letter to Quillen by distinguishing the study of spaces of maps and their successive homotopies from the study of spaces of homeomorphisms, immersions, embeddings and fibrations. His terms ``wildness'', ``isotopy'', ``stratified structures'' and ``tubular neighbourhoods'' delimit problems for which the familiar category of topological spaces seems poorly adapted; the contrast with piecewise linear structures is that their maps are too sparse, notably because real-valued piecewise linear functions are not closed under multiplication. This develops the complaint recalled explicitly from the letter, p.~1, and supplies the contrast on which section~15 builds its more favorable account of localization and homotopy types. The interlinear revisions narrow the opening description to ``maps between spaces'' and replace ``numerical'' by $\mathbb R$-valued (p.~2). The witness differs in two mathematically relevant places: the typescript says ``pro-topoi'' on p.~1, where arXiv has ``pretopoi'', and arXiv retains ``numerical'' without applying its handwritten replacement on p.~2. These readings are established by the scan; the edition does not identify the two notions or silently repair either authorial formulation.

\subsection*{15. Models of homotopy types and the place of stacks (pp. 2--4)}
The argument passes from the failure of the Hauptvermutung as a literal assertion about triangulations to its much weaker homotopical analogue: appropriate localizations of topological spaces and semi-simplicial sets yield the same homotopy category. Grothendieck compares CW-spaces, Kan complexes and Serre's cubical complexes, then proposes stacks, previously called ``$\infty$-groupoïdes'', as a third description particularly close to the homotopy groups themselves. His analogy with chain complexes and the Dold--Puppe equivalence makes the intended role of noncommutative structure explicit; section~16 immediately generalizes the motivation from spaces to topoi. The distinction between describing and computing invariants matters: on p.~3 he twice strikes ``description'' and writes ``computation'', corrections which the arXiv witness does not apply. On p.~4 a marginal note dated 3.5 begins ``This is an error'' and retracts the claim that stack cohomology has no immediate description from the primitive structure, referring instead to the nerve and to forgetting structure. Broad pencil strokes across this note make its final cancellation status uncertain, and one verb remains doubtful; both uncertainties are preserved. The arXiv footnote is substantially longer than this marginal wording and must not be substituted for it as though it were a literal transcription.

\subsection*{16. Stacks over topoi and noncommutative cohomological algebra (pp. 4--6)}
The motivation for stacks is enlarged from section~15's models of homotopy types to coefficients for cohomology over an arbitrary topos. Starting with Giraud's $0$-, $1$- and $2$-dimensional theory, Grothendieck treats $n$-stacks as coefficients whose ``integration'' is evaluation at the final object, with the homotopy groups of the resulting stack supplying cohomology after a dimension shift. His letters to Breen and Mme Raynaud's work on Lefschetz comparisons provide the geometric motivation: higher noncommutative cohomology should unify homological and homotopical comparison statements. The passage to infinite-order stacks is presented as a qualitative enlargement of perspective, and the following section asks how much of the homotopy-theoretic side ordinary categories can already model. The scans matter for the formulas: p.~4 has $H^i$, not the witness's $K^i$; p.~5 has $R\Gamma F_*$ and complexes $F_*$, where the supplied arXiv comparison reports a missing $F$ and a $j$ in place of the star. The handwritten qualification ``[non commutative]'' on p.~5 must also remain visible. These are witness differences, not authorial mathematical corrections introduced by this edition.

\subsection*{16$'$. Categories and the ``impressive bunch'' of models (pp. 7--9)}
Grothendieck associates a homotopy type to a small category both through its presheaf topos and through its nerve, comparing the first construction with Verdier and Artin--Mazur pro-homotopy types. He recalls and rejects his former guess that passage to the enveloping groupoid preserves this type: Quillen had pointed out that ordinary categories realize arbitrary homotopy types. This corrects the initial groupoid intuition while continuing sections~15--16's search for models, and motivates the later study of the ``impressive bunch'' and of coherators. A second proposed shortcut is itself explicitly retracted: the margin of p.~7, dated 5.3, says that the asserted weak equivalence $K\longrightarrow NN'(K)$ is false and directs the reader to section~24, pp.~21--22. The note is shorter than its arXiv counterpart; its call number is (18), and the section has been renumbered $16'$ rather than left as $16$bis. Page~8 proposes cofibrations injective on objects and arrows with a further call (19), then contrasts general $n$-stacks with $n$-Gr-stacks and their truncated homotopy types. Its wording ``left derived functors'' for the inverse-limit cohomology construction should be treated as an authorial formulation requiring care, not silently changed to the usual right-derived terminology. The final paragraph on p.~9 introduces the dependence on a coherator and the universal-algebra question, despite professing little inclination to pursue it.

\subsection*{17. The Artin--Mazur cohomological criterion (p. 9)}
This short section supplements the unification program by recalling that weak equivalences can be detected through cohomology with twisted coefficients, including noncommutative $H^0$ and $H^1$ as well as the commutative higher groups. Grothendieck explains how this permits étale cohomological information to constrain homotopy types, while wondering whether its success discouraged the development of higher noncommutative cohomology. The reference back to section~15 is the author's own handwritten cross-reference and is retained even though the immediate unification discussion occupies section~16 in the present numbering. The section bridges that motivation and the concrete corrections of section~18. Its handwritten insertion reads ``(namely, unification)'', not the arXiv ``messy unification''; the final sentence deletes ``that'' and says ``their observation'', distinctions that are obscured by the witness. The section number itself is corrected from 18 to 17 on the scan.


\subsection*{18. Corrections to the coherator and higher coherence (pp. 10--11)}
The return to the letter (L) is explicitly corrective. Grothendieck withdraws the existence of arbitrary binary amalgamated sums in $\underline B_o$: the geometric intuition of gluing contractible cells along embedded subcells had concealed the presence of nonmonomorphic maps, including the codiagonal. He also revises the construction of $\underline C_{n+1}$ from $\underline C_n$ by insisting on the total set of new arrows, and rejects the assertion that intersections of maximal subcells must be subcells. The marginal drawing on p.~10 makes the latter failure concrete: three regions $f,f',f''$ meet so that $f\cap f''=\{s,t\}$. The scan therefore supplies a two-point intersection, not the quite different formula in the initial transcription. A separate correction explains why associativity beyond primary composition cannot be represented by a single homotopy arrow: the sources and targets are themselves only homotopic, so squares, cubes and higher coherence data intervene. Page~11 names the dual categories left and right coherators, then connects the proposed non-groupoidal stacks to Benabou's bicategories and their cellular, topological interpretation. This prepares section~19's realization and singular-stack adjunction. The source wording $L\to K$ and $M\to K$ beside $K\amalg_L M$ on p.~10 is retained as an apparent authorial slip; it does not give the usual span for that pushout. Calls (11) and (14) on p.~11 are present without the explanatory footnotes printed in arXiv.

\subsection*{19. A provisional itinerary and the sheaf temptation (pp. 12--13)}
The first two items of the itinerary are the definition of $\underline B_o$, its dual and the coherators, followed by a realization/singular-stack adjunction with topological spaces. Unlike arbitrary presheaves on a left coherator, Gr-stacks must carry the designated amalgamated sums to fiber products. The realization functor nevertheless extends to all presheaves by preserving colimits; its values on basic cells are prescribed, while extensions to the coherator require choices. The anticipated independence of those choices is therefore tied to an as-yet undeveloped notion of equivalence between Gr-stacks. On p.~13 Grothendieck asks whether the exactness conditions could instead be sheaf conditions for a topology, which would make Gr-stacks a topos, and proposes weak equivalences detected by realization, monomorphisms as cofibrations and a lifting-property definition of fibrations. These are programmatic proposals rather than established results. The ``quick glance'' suggesting a topology is retracted immediately in section~20; even the concluding sentence about the inclusion, revised to composite, functor is already struck on p.~13. The source has ``one reward'', not the witness's ``one ward'', and handwritten ``as cofibrations'', not ``for cofibrations''. Its call (13) contains no footnote text on this page.

\subsection*{20. Model categories, sites and a retracted canonical topology (pp. 14--15)}
The section opens by correcting section~19's ``quick glance'': the proposed topology on the coherator would make the component functors fiber functors and force componentwise colimits, which the author recognizes as false. The category of categories supplies a warning and a new candidate construction. Grothendieck considers families of functors whose nerves cover every finite composable string, proposes identifying them with universally effectively epimorphic families, and associates to a category $A$ the topos of sheaves on $(\mathrm{Cat})_{/A}$. The handwritten margin immediately marks the asserted equivalence as false and points to section~24; the digits are readable, while pencil hatching leaves the note's cancellation status uncertain. Page~15 broadens the proposal to a natural site structure on model categories, with the homotopy type of $X$ identified with that of the induced site or topos over $X$. Its further presumption that the topology should be canonical is itself contradicted by a dated margin note absent from the arXiv comparison. These local corrections are essential to reading the section as an evolving investigation. The four arrows in the diagram of effective, universal and ordinary epimorphisms are handwritten, as are the topos tildes in $(M_{/X})^{\sim}=M^{\sim}_{/X}$; the object labels are typed. The inquiry leads into section~21's wider inventory of models rather than establishing a general theorem about canonical sites.

\subsection*{21. Multicomplexes, symmetry and restricted coherence (pp. 15--16)}
The inventory now includes symmetric variants of semi-simplicial and cubical complexes and their multiple-complex analogues. A diagonal complex already describes the intended homotopy type, but Grothendieck explicitly pursues the variety of algebraic presentations, not a unique most convenient model. He compares the passage from complexes to multicomplexes with that from ordinary categories to higher categories or stacks, recalling Quillen's suggestion of defining higher categories through multiple semi-simplicial complexes. The return to the coherator is concrete: the coherence diagrams examined in Benabou's bicategory axioms use restricted amalgamated sums describable through multisimplices, suggesting a smaller starting category $\underline B_o$ than the generous one in the original letter. The handwritten cross-references on p.~16 revise the old page references to sections~8 and~9, and the scan preserves the tense and spelling ``This I ckecked'', where arXiv reads ``This I check''. The suggested nerve-like functor from $n$-stacks to $n$-semi-simplicial complexes is a goal for later investigation. Section~22 next reverses direction, seeking simpler special cases of categories, beginning with preorders and ordered sets. The uncertain final singular/plural in ``connection(s)'' on p.~15 remains an image-edge problem and has no mathematical consequence for this comparison.

\subsection*{22. Ordered sets, realizations and maquettes (pp. 16--19)}
Preorders and ordered sets offer simpler categorical presentations of homotopy types. He relates ordered strata in conically stratified spaces to a realization whose barycentric simplices are finite totally ordered subsets, or flags. The resulting realization is canonically compared with that of the nerve. On p.~17 the corrected chain is $(\mathrm{Ord})\hookrightarrow(\mathrm{Preord})\hookrightarrow(\mathrm{Cat})\hookrightarrow(\mathrm{Ss\ sets})$, not the reversed first pair in arXiv. He cautions that the adjoint passage from a category to its object preorder can lose homotopy information, as the one-object case shows; the handwritten qualification ``but not Fl(A)'' is indispensable to that example. Page~18 gives a second realization through a nonseparated topology on the underlying ordered set and identifies sheaves with covariant functors, hence presheaves on the opposite order. Reversing the order preserves the flags and yields homeomorphic geometric realizations, while the two associated presheaf topoi need not be equivalent. The scan confirms hats in the latter claim and the author's rounded inclusion sign for the initial order comparison. Finally, maquettes are described by vertices and finite simplices and embedded in ordered sets through their face posets. The body claims that this category lacks finite products, while a dated marginal note adds that the notion is not an algebraic structure species; the hatching and call number of that note remain uncertain. Section~23 consequently shifts attention from satisfying Quillen's full package to constructing and localizing a natural functor into $(\mathrm{Hot})$. The edition preserves these claims and reservations as stages of the investigation, without treating the later witness's explanatory footnotes as text written on these pages.

\subsection*{23. The functor to homotopy types and a cellular core (pp. 19--21)}
Grothendieck separates the initial construction of a natural functor $M\to(\mathrm{Hot})$ from the further demand that $M$ carry a Quillen model structure. The functor would define weak equivalences by inverse image of isomorphisms; only then would one choose monomorphisms as cofibrations and test the resulting lifting-property definition of fibrations. The proposed site-theoretic construction uses the homotopy type of the induced site over each object, provided its pro-homotopy type is essentially constant. The canonical-topology suggestion is already qualified by a handwritten retraction pointing to section~24. Page~20 adds a structural requirement, universal and disjoint coproducts, and presents an algebraic species by a category $\underline C$ with finite inverse limits, its dual $\underline B$, and the anticipated equivalence $\mathrm{Ind}(\underline B)\simeq M$. This connects localization to finite presentations and to the question of which exactness properties are already visible in $\underline B$. Page~21 then isolates a smaller category $\underline B_{\mathrm{oo}}$ of basic cells and monomorphic boundary maps: the ordered set of proper subobjects should realize the boundary sphere of the corresponding cell. That construction uses section~22's ordered-set realizations and develops section~18's correction of the initial gluing intuition. The distinction between the broad species-level proposal and the geometric cell data remains heuristic. In the scan, the long p.~20 insertion includes ``can be interpreted'', omitted by arXiv, and the identity-functor arrow is drawn with two shafts; p.~21 preserves the printed slip ``of cubical'' and the typed replacements ``define'', ``family'' and ``sphere''.

\subsection*{24. Truncated nerve topologies and the failure of uniqueness (pp. 21--24)}
This section supplies the correction repeatedly anticipated in the earlier margins. Covering all composable pairs does not cover arbitrary strings: Grothendieck defines a descending sequence $T_N$ by epimorphic families of $N$-truncated nerves, identifying the first two levels with universal epimorphy and universal effective epimorphy. The comparison lemma relates the associated sheaf topoi to presheaves on truncated simplex categories. The finite levels retain only skeletal information, so the full nerve topology $T_\infty$ is needed for the proposed homotopy-type construction. The claim on p.~22 that an $N$-skeleton has the same homotopy and cohomology invariants through dimension $N$ requires caution at the endpoint: the $N$-skeleton of an $(N+1)$-simplex is an $N$-sphere, although the full simplex is contractible. The edition preserves the author's inequality rather than silently strengthening or repairing this intermediate assertion. The same loss of information explains the failure of $K\to NN'(K)$ discussed in section~16$'$: $N'$ factors through the two-dimensional data, while the nerve itself can still induce the desired equivalence after localization. Page~23 then replaces simplicial diagrams by commutative cubical diagrams, producing a strictly coarser topology with a different associated topos but the same intended homotopy interpretation. Multiple complexes suggest infinitely many further suitable topologies; the conclusion is that no unique topology is privileged. This leads directly to section~25's search for an equivalence relation among them. The scan's slice notation on p.~22 places $/A$ outside the parentheses ending after $T_N$, followed by an additional handwritten closing parenthesis; the earlier transcription's placement inside the subscript was corrected, while the extra parenthesis remains as written.

\subsection*{25. Hot-equivalence of topologies and global asphericity (pp. 24--25)}
The lack of a privileged suitable topology leads to a proposed equivalence relation. For comparable topologies $T\geq T'$ on $M$, Grothendieck compares the induced morphism of sheaf topoi and asks that every induced morphism over an object $A$ preserve homotopy type. He calls the resulting property ``globally aspheric'', strengthening the familiar cohomological notion of global acyclicity by allowing noncommutative coefficients. The primitive relation between comparable topologies is not yet an equivalence relation; taking the relation it generates is a first proposal, and he asks whether a coarser useful notion could do better. The intended invariant is a canonically associated functor $M\to(\mathrm{Hot})$, perhaps strong enough to recover the topology class itself. The analogy with Giraud's recovery of a topology from its associated subtopos motivates, rather than proves, a hoped-for correspondence. Page~25 narrows the aspiration to a single natural Hot-equivalence class for suitable species, admitting that even $M=(\mathrm{Cat})$ has no intrinsic characterization yet. Section~26 then reformulates the problem through small subcategories. Two revisions on p.~24 are essential: ``inverse image'' becomes ``direct image'', and the prime is moved so that the inclusion is from $T$-sheaves into $T'$-sheaves. The vertical margin note actually says ``just means'' and ``namely contractible'' with the qualification ``in the Čech sense''; the initial transcription's ``weakly contractible'' was not a literal reading. Its pencil shading leaves cancellation status uncertain. The body ends with the typed slip ``he'' for ``be'', and its final ``this'' is visibly canceled before the continuation on p.~25.

\subsection*{26. Test categories, test functors and the passage back to $M$ (pp. 25--29)}
The proposed intrinsic characterization of suitable topologies is reformulated through small categories $B$ inside a model category $M$. The cellular cores of section~23 and the simplex and cube categories of section~24 suggest that the presheaf category $\widehat B$ should itself describe homotopy types: its natural functor to $(\mathrm{Hot})$, computed from categories of objects over a presheaf, should be a localization. Grothendieck separates this internal requirement on $B$ from the eligibility of a functor $B\to M$; marginal terminology on p.~27 calls them ``test categories'' and ``test functors''. The comparison lemma gives an equivalence between the induced sheaf topos and $\widehat B$ under either of two fullness hypotheses. The decisive qualification comes on p.~28: localization of that sheaf category is not yet localization of $M$ itself. Generation by strict epimorphisms is considered as a way of making $M\to\widehat B$ fully faithful, but the author then proposes dropping both fullness requirements and directly demanding that $M\to\widetilde M\to\widehat B\to(\mathrm{Hot})$ be a localization. For categories $\underline B$ of finite presentations he further hopes for a localization onto homotopy types of finite type. The questions whether equal weak equivalences determine the functor, and whether suitable choices form one natural class, feed sections~27--28. The margin of p.~27 actually says ``same question arises'' for Hot-equivalence of topologies; ``arises'' was absent from the initial transcription. Its wording is now readable, but pencil hatching still leaves the note's cancellation status uncertain. The call after objects ``over a'' on p.~26 is (23), not (3).

\subsection*{27. Geometric realization as a third construction (p. 29)}
Calling the preceding discussion ``thin air conjecturing'', Grothendieck warns that both test categories and test functors may need stronger conditions. He places geometric or ``topologic'' realization beside the two constructions through topologies and through test functors. Semi-simplicial and cubical complexes offer comparatively direct realizations, hence so do the categories embedded in semi-simplicial sets. For Gr-stacks, the prescribed realization of the primitive $\underline B_{\mathrm{oo}}$ structure has to extend to a left coherator, and that extension involves the arbitrary choices already encountered in sections~18--19. Rather than immediately eliminating them, he compares this indeterminacy with the variety of test topologies and functors developed in sections~24--26. The paragraph thus states the next intended problem, relating realization to the two test constructions; section~28 immediately postpones it in favor of rigidity of $(\mathrm{Hot})$. The scan preserves ``topologic'', where the witness expands it to ``topological'', and the word ``is'' in ``is that a suitable'', where the witness has ``it''. Its call (28) supplies no text of the witness's retrospective footnote. The typed insertions ``localization'', ``left'' and ``between'' sharpen the objects of the comparison; the two round letters in $\underline B_{\mathrm{oo}}$ and the primitive graph designation remain literal rather than silently replaced by a joined infinity glyph.

\subsection*{28. Rigidity of $(\mathrm{Hot})$ and the first modelizers (pp. 30--33)}
The intended return to geometric realization is suspended by an ``Assumption'': the category of self-equivalences of $(\mathrm{Hot})$ and their natural isomorphisms should be equivalent to a point. Grothendieck separates its two demands, every self-equivalence is isomorphic to the identity and the identity has no nontrivial automorphism, while wondering whether arbitrary natural endomorphisms could also be admitted. This is asserted as a plausible assumption, not proved; it supplies the desired uniqueness up to unique isomorphism of a localization functor once its weak equivalences are known. A ``modelizing category'' or ``modelizer'' is then a category $M$ with a saturated class $W$ whose localization is equivalent to $(\mathrm{Hot})$. A model-preserving functor carries weak equivalences to weak equivalences and induces an equivalence of localizations, equivalently commuting up to isomorphism with the type functors. This formalizes sections~26--27's questions without yet specifying cofibrations and fibrations. The diagram on p.~31 compares ordered sets, preorders, categories, semi-simplicial sets and cubical sets; the category-of-elements functors $i_A$ lead to right adjoints $j_A$. Page~32 carefully distinguishes those right adjoints from the familiar nerve functors, contrasting infinite categories of simplices with finite ordered face structures when degeneracies are omitted. Mixed multicomplexes then enlarge the diagram and motivate the test-category definition in section~29. The mixed-product display on p.~32 visibly has $\Delta$ in both factors, although the adjacent prose says $p$ times simplicial and $q$ times cubical; this apparent authorial symbol slip is preserved. The same page changes ``modelators'' to ``modelizers'' by hand, a correction the witness does not apply, and contains a small branching diagram omitted there. The p.~30 cross-reference is to section~26; calls (28$'$), (29) and (30) are present without the witness's expanded footnotes.

\subsection*{29. The basic modelizer and a stronger test condition (pp. 33--36)}
The category of small categories now becomes the foundational modelizer: $(\mathrm{Hot})$ is its localization, and weak equivalences are described through the induced morphisms of presheaf topoi and the Artin--Mazur cohomological criterion recalled from section~17. A small category $A$ is provisionally a test category when $i_A\colon\widehat A\to(\mathrm{Cat})$, sending a presheaf to its category of elements, is model preserving; its presheaf category is an ``elementary modelizer''. This replaces a special preference for semi-simplicial sets by a comparison among all such presentations, while acknowledging the computational convenience of the nerve and the simplicial construction of pro-homotopy types. The general adjunction calculation of pp.~34--35 asks when both adjoints induce inverse equivalences after localization. Applied to $i_A\dashv j_A$, it leads to condition (T): every counit $i_Aj_A(B)\to B$ should be a weak equivalence. Grothendieck explicitly leaves open whether this strengthens the preceding definition, and proposes strengthening it if necessary; the call (33) does not contain the witness's retrospective answer. Products of test categories and direct comparison of elementary models become the next questions. The closing heuristic on p.~36 would test (T) by a fibration with contractible fibers, either after taking nerves or directly in $(\mathrm{Cat})$; section~30 immediately attacks the model-category assumption behind the latter route. On p.~35 the source twice writes that $M'$ or $M$ is saturated, although the discussion concerns the arrow classes $W'$ and $W$; these apparent symbol slips remain unchanged. The typed ``elementary normalizer'' on that page likewise remains source text. The scan's baseline slash in $M/_{b}$ on p.~34 is preserved, and its reference reads ``section 28'', not the witness's page reference. The broad emphasis of the witness also differs from the selective underlining and marginal double bar of p.~33.

\subsection*{30. From the model-category obstruction to asphericity (pp. 36--38)}
The attempt to treat every monomorphism of categories as a cofibration now runs into the two elementary horns drawn on p.~37. Lifting against both would force the arrows of a fibrant category to have left and right inverses, leaving groupoids, whose homotopy types cannot exhaust those of arbitrary categories. The objection concerns that proposed choice of cofibrations; the source does not prove the nonexistence of every possible model structure on categories. Grothendieck nevertheless keeps the expectation that elementary modelizers should admit closed model structures, and turns instead to the asphericity criterion introduced in section~25. A morphism of topoi is required to remain a weak equivalence after every localization over the target; when the topoi come from categories, generators reduce the condition to the asphericity of each slice $C'_{/a}$. Applying this to the counit $i_Aj_A(B)\to B$ from section~29 gives three equivalent conditions: the strengthened test condition, asphericity of every counit, and asphericity of $i_Aj_A(B)$ whenever $B$ has a final element. That last formulation opens the explicit interval argument of section~31. The scan's four bracketed ``[small]'' restrictions on p.~38 are genuine hand additions, including two absent from the witness. Its caret under the $A$ in $\operatorname{Hom}(A_{/a},B)$ is retained in its actual position without asserting a mathematical interpretation. The apparent inline hat on $f$ reported by the initial comparison on p.~37 belongs instead to the display beneath it; the inline symbol has been corrected to plain $f$. The word ``closed'' on p.~37 and the replacement ``aspheric'' on p.~38 are interlinear typing, although their placing strokes are handwritten.

\subsection*{31. An interval proof of the test-category criterion (pp. 38--41)}
The search for a direct proof that the simplex category is a test category yields three sufficient conditions. T1 asks that $A$ be aspheric; T2 asks that $A_{/a\mathrm{x}b}$ be aspheric for any two objects, with their product interpreted in the presheaf topos when necessary; T3 supplies an aspheric interval object $I$ with disjoint final subobjects $e_o,e_1$. For $F=j_A(B)$ with $B$ admitting a final element, the aim is a homotopy from the identity to the constant section. Diagram (D) on p.~39 expresses the cylinder $F\mathrm{x}I$ and its two endpoint inclusions. Given $u:a\to I$, the inverse image of $e_o$ determines a sieve $C_u$ in $A_{/a}$; a functor $f$ is extended by its own values on $C_u$ and by the final object of $B$ on the complementary objects. The sieve property makes the remaining arrows uniquely definable. This constructs the homotopy without T1, T2 or the asphericity of $I$. Those hypotheses enter the second step: the slice criterion of section~30 reduces the asphericity of $F\mathrm{x}I\to F$ first to $a\mathrm{x}I$ and then to $a\mathrm{x}b$, precisely T2. T1 then converts the weak equivalence to the final presheaf into absolute asphericity. Page~40 separately justifies the retract step through the Artin--Mazur criterion rather than claiming it is formal for an arbitrary arrow class. The scan gives genuinely hooked tails to both arrows into the cylinder in (D), contrary to the initial transcription; all diagram arrows are handwritten and its object strings are typed. The p.~39 reference is ``cf théorème p 46'', while the marginal comparison involving T1 and T2 retains an uncertain short connector. The live punctuation after ``checked'' on p.~41 is a colon, and the source's ``asperic'' remains uncorrected.

\subsection*{32. A revised programme after the interval criterion (p. 42)}
Dated 15.3, the new planning page treats T1--T3 as convincing candidates rather than a finished characterization. Its first three tasks are to examine necessity, construct the familiar simplicial and cubical examples, and prove stability under products so that different elementary models can be compared. The fourth returns to section~30's unresolved closed-model question; the fifth restores the larger programme of modelizers attached to algebraic structure species and the possible uniqueness of their weak equivalences. The author explicitly notices that stacks have receded into the background. A sixth task broadens the inquiry to a satisfactory replacement for Verdier's provisional triangulated-category notion, including the exactness, generator and cogenerator questions for $(\mathrm{Hot})$. Section~33 begins carrying out the first task and introduces the distinction between test and strict test categories. The scan matters to the direction of the programme: item 5 says categories $M$, not the witness's $\widehat A$. Its call (9) is present without the witness's explanatory footnote. ``Gave to this notion'' and ``general'' are typed interlinear revisions, while ``ways of'' and the displaced ``of'' are handwritten. The cancelled ``topos'' in item 4 is legible beneath both machine and pen deletion; the cropped completion of ``which'' and the old word after ``constructing'' remain doubtful.

\subsection*{33. Strict test categories, the Lawvere interval and a qualified theorem (pp. 43--47)}
The provisional terminology first separates test categories from ``strict'' test categories satisfying T1--T3. T1 is necessary by the final-category case of the criterion on p.~38. T3 admits a canonical candidate: applying $j_A$ to the two-object ordered interval produces a presheaf whose sections are sieves, hence the ``Lawvere element'' representing subobjects. This ties the explicit interval construction of section~31 to an intrinsic object of the presheaf topos. A proposed concluding statement about an intrinsic T3$'$ is completely cancelled on p.~43. The remaining issue is T2. Pages~44--45 connect it with preservation of binary products by the localization functor: the comparison $A_{/F\mathrm{x}G}\to A_{/F}\mathrm{x}A_{/G}$ is a weak equivalence for all $F,G$ exactly when the representable-product condition holds, and the slice criterion in fact supplies asphericity. The blanket claim about finite direct and inverse limits in $(\mathrm{Hot})$ on p.~44 is stronger than the finite sums and products used here; its marginal (!35) must not be expanded into an unqualified theorem in the guide. The result assembled on p.~46 compares product-preserving localization, the localized adjunction $i_A\dashv j_A$, and T1--T3. Its later margin dated 9.4 is essential: condition (i) is said to be ``not strong enough'', and the left continuation proposes adding that $j_A=i_A^*$ is model preserving, with a reference to remark~3 on p.~174. The left note is readable but crossed by broad pencil hatching, so its cancellation status remains uncertain. The source says ``model preserving'', not the witness's ``equally modelizing''. The author also explains how to avoid assuming that $(\mathrm{Cat})\to(\mathrm{Hot})$ preserves binary products: formulate the condition directly for $i_A$ up to weak equivalence. Page~47 provisionally reserves ``test category'' for the strict case, with call (38), and characterizes elementary modelizers through presheaf-topos structure, localization and finite products; a wider modelizing-topos notion is deferred. The scan retains ``asphéric'' on p.~43, three rounded intersection caps rather than the transcription's wedges, typed interlinear additions to the theorem, and the actual oval marginal annotations. The two opposite bent arrows on p.~47 remain graphic annotations, without an invented reordering instruction.

\subsection*{34. Product-stable examples and the homotopy circularity (pp. 47--50)}
The first examples strengthen the sufficient hypotheses again. A final object gives T1$'$, and an actual object $I$ of $A$ with two disjoint sections gives T3$'$; if $A$ is also stable under binary products, T2 follows because each product is representable and its slice has a final object. This constructs test categories inside any suitable ambient category, even one that is not small, by selecting a small product-stable subcategory. The economical construction uses powers of an interval and only projections and constant endpoint maps as components. Its cubical category includes permutations and diagonals, so p.~49 distinguishes it from the more familiar semicubical category. Nonempty finite sets give a corresponding full simplicial construction, with affine simplices as its geometric interpretation. The notation on p.~48 literally says to include $I$ and $\Delta_o,\Delta_1$; these triangular symbols are retained even though the preceding discussion specifies endpoint maps $\delta_o,\delta_1$. Returning to the ordinary simplex and cube categories removes the convenient product-closure hypothesis. Grothendieck therefore tries to prove T2 by making every representable deform to the final presheaf. The decisive revision on p.~50 is his own detection of a circular argument: section~31's homotopy proof used $F\mathrm{x}I\to F$ as a weak equivalence, and that was established from T2. A separate condition, called (2), that every $a\mathrm{x}I$ be aspheric must therefore accompany the deformation-retract condition TH1. Section~35 postpones this unfinished verification for the intrinsic theory of modelizing topoi. The formula $u_a=p_a c_a$ on p.~50 is retained in that visible order; with $c_a:e\to a$ and $p_a:a\to e$, the surrounding right-to-left composition convention would instead require $c_a p_a$ for an endomorphism of $a$. The scan also confirms the literal spellings ``axists'', ``oddurs'' and ``consequance''. The extra presheaf caret on the full simplicial category on p.~49 is present despite the witness omission, and the stepped mark on p.~48 is retained graphically in the diplomatic reading and applied as a paragraph break only in the master reading.
\subsection*{35. Local and total asphericity of topoi (pp. 50--55)}
The detour from section~34's unfinished simplex proof translates its conditions into the language of topoi. An aspheric topos is required to be $0$-connected and to have trivial cohomology with constant coefficients in every defined positive degree, including the noncommutative first degree. For an object or arrow, the induced topos supplies the corresponding notion. A locally aspheric topos has a generating family of aspheric objects; these generators suffice for testing whether an arrow remains a weak equivalence after every localization (p.~51). Proposition~1 identifies product stability of aspheric objects with their being aspheric over the final object. Total asphericity combines that property with local asphericity, and the interlinear ``non empty'' inserted on p.~52 is essential: without it, the preceding assertion that T2 already implies T1 would admit the empty case. The ensuing \v{C}ech argument gives absolute asphericity. For topological spaces the much stronger total condition becomes irreducibility, so a Hausdorff example must be a point. This sharpens section~33's distinction between a local supply of cells and the global product condition.

Proposition~2 then separates the interval requirement from total asphericity. Its original ``totally aspheric'' hypothesis is visibly cancelled, and p.~53 explicitly says that no local-asphericity assumption is being made. The remaining equivalence is between the Lawvere element being aspheric over the final object and the existence of an object with that property and two disjoint sections. The homotopy argument of section~31 supplies the converse, and localization upgrades the resulting weak equivalence to asphericity. A connected topological space provides a useful obstruction: the proof with direct summands of its Lawvere element shows that this element is disconnected precisely in the irreducible case, where it has two connected components. Consequently no topological space can simultaneously meet total asphericity and the interval requirement (pp.~53--54). The proposed term ``modelizing topos'' packages both, while a second, still conjectural characterization asks for a saturated localization equivalent to $(\mathrm{Hot})$, compatible with finite sums and products and with the Verdier--Artin--Mazur construction. The closing inquiry about morphisms $X\to\operatorname{Top}(B)$ seeks a categorical description closer to the given indexing category than a large inverse system of semi-simplicial objects (pp.~54--55). The source retains ``0-connexity'', ``independant'', ``Corollay'' and ``monomorphisme''. Its upper margin on p.~54 reads ``cf (34)'' beside a body call (39); the lower strictness-related note remains partly doubtful. These are retained source features, not expanded into the arXiv witness's retrospective commentary.

\subsection*{36. Shuffles, nerves and comparison with categories (pp. 55--58)}
The simplex proof postponed in section~34 is completed by separating its two tasks. The ordered map $\Delta_n\mathrm{x}\Delta_1\to\Delta_n$ contracts toward the last element, giving TH1. For TH2, the maximal flags of the product order---the ``shuffles''---are attached successively so that each new intersection is a subflag. A Mayer--Vietoris-type lemma then propagates asphericity from two subobjects and their intersection to their union (p.~55). The author's inserted ``and covering e'' is indispensable to that lemma. He immediately warns against replacing this argument by the observation that the product ordered category has a final element: that shortcut assumes the still-unproved compatibility between the category's homotopy type and that of its nerve. The literal typed replacement ``apshericity'' remains uncorrected in the reader. The adjacent handwritten call is (41), not the asterisk proposed by the first transcription.

The comparison proposition of pp.~56--57 addresses the missing compatibility in a more general form. A functor $a\mapsto\Delta_a$ from $A$ to categories defines $\alpha(C)(a)=\operatorname{Hom}(\Delta_a,C)$. If each $\Delta_a$ has a final element, $A$ is aspheric, and $\alpha(\Delta_1)$ is aspheric over the final presheaf, choosing final elements defines an aspheric functor $\varphi:A_{/C}\to C$, natural in $C$. The choice of final rather than initial elements makes the extension to arrows possible. Its slices identify with the corresponding construction over $C_{/x}$, so it suffices to treat $C$ with a final element; applying $\alpha$ to the ordinary categorical contraction supplies the required interval homotopy. Thus $i_A\alpha$ preserves weak equivalences and induces the identity on $(\mathrm{Hot})$. The simplex case now needs only the verified TH2, while the category of nonordered simplices is sent to the ordered set of nonempty faces of each finite set. A proposed sieve criterion for the interval condition follows on p.~58, but the entire corollary and its proof are diagonally cancelled before the new date 21.3. Its uncancelled presentation in arXiv must not conceal that revision. The source diagram on p.~56 literally labels its upper objects $a',a$ rather than $\Delta_{a'},\Delta_a$; the reader preserves this abbreviation or slip. The cancelled formula on p.~58 has $C_u\cap C_o$, not the witness's composition sign, and its first $F\to e$ occurrence is actually drawn with a mapsto tail. Section~37 turns the repeated argument into a new systematic account.


\subsection*{37. Homotopy intervals and the key comparison theorem (pp. 58--62)}
The new account dated 21.3 begins by identifying the gap in section~36's cancelled sieve criterion: the choices $C_u$ must be functorial in $u$, so that they actually define a map to the proposed interval. The replacement of ``automatic'' by a claim restricted to the intended examples makes that qualification explicit (p.~59). Grothendieck now formulates the argument in a category with finite products and a class $W$ satisfying isomorphism, two-out-of-three and retraction conditions. The last condition is stronger than the earlier mild saturation convention. An interval has two disjoint sections and an aspheric projection, with the required base changes understood; the margin calls the terminology ``provisional'' and adds ``separating''. The first homotopy lemma contracts an object to a section and concludes universal asphericity, while explicitly observing that disjointness was not used. The comparison lemma then transfers this property along an endpoint-preserving map to an object equipped with a unit-and-zero multiplication (p.~60). The revised formulas consistently make endpoint 1 the unit and endpoint 0 the zero, although the earlier ``respectively'' sentence is not fully adjusted. The Lawvere element supplies the universal example. In its proof on p.~61, the pullback along $e_0\to L$ is still called a subobject of $e_1$; the edition preserves this apparent index slip.

The resulting theorem separates two strengths of comparison for a functor $i:A\to(\mathrm{Cat})$ whose values have chosen final objects. The choices need not be preserved by the arrows of $A$. Part (a) equates asphericity of the evaluation $A_{/C}\to C$, its being a weak equivalence, and the resulting identity on localized homotopy types, with equivalent tests on categories having a final element. Part (b) asks for asphericity relative to $A\times C$, equivalently for the appropriate representable products, and reduces this stronger requirement to universal asphericity of $i^*(\Delta_1)$. Together with asphericity of $A$, part (b) implies part (a); the first remark explicitly doubts the converse (pp.~61--62). The proof again transports the canonical categorical contraction through the left-exact functor $i^*$, while nonemptiness of each $i(a)$ ensures separation of the two endpoints. A superseded proposed condition in part (b) is cancelled by both machine and pen. The marginal reference to the sequel in volume~2 has a still-uncertain first page numeral, tentatively 64--65. Section~38 extracts the consequences for the Lawvere criterion and asks how far the final-object hypothesis can be removed.

\subsection*{38. The Lawvere criterion and generalized nerves (pp. 63--66)}
The theorem's main gain is the reduction to one interval condition: universal asphericity of $i^*(\Delta_1)$, together with asphericity of $A$, makes $i_Ai^*$ induce the identity on $(\mathrm{Hot})$. Both that interval condition and the localized conclusion make sense without chosen final objects in the categories $i(a)$, and p.~63 therefore anticipates a more general comparison theorem, while acknowledging that a direct evaluation map would then be unavailable. For the canonical choice $a\mapsto A_{/a}$, the right adjoint is $j_A$ and $j_A(\Delta_1)$ is the Lawvere element $L_A$. Corollary~1 consequently gives a sufficient criterion for a test category in the wide sense: $L_A$ is aspheric over the final presheaf and $A$ is aspheric (p.~64). The scan actually has a separate handwritten caret in ``unit element \ldots\ of $\widehat A$'' on p.~63; the initial transcription's bare $A$ was not a source slip.

For general $i$, a crible in $i(a)$ determines the crible of arrows $b\to a$ whose images factor through it. This defines an endpoint-preserving map $i^*(\Delta_1)\to L_A$, compatible with intersection. Section~37's comparison lemma transfers the interval condition to $L_A$. Corollary~2 then makes both $i^*$ and $i_A$ modelizing and identifies the functor induced by $i^*$ with the canonical quasi-inverse induced by $j_A$ (pp.~64--65). Thus the localized functor does not depend on the admissible choice of $i$, although the resulting concrete presheaf models can be much simpler than those obtained from the often infinite slice categories $A_{/a}$. The ordinary nerve is the motivating example of this ``generalized nerve functor''. Corollary~3 makes the remaining verification practical: if $A$ has a final object $e$, $i(e)=\Delta_0$, and a homotopy interval $I$ admits an endpoint-preserving map $i_!(I)\to\Delta_1$, adjunction and the same comparison lemma give the desired asphericity of $i^*(\Delta_1)$ (p.~66). In the intended examples $I$ belongs to $A$ and $i(I)=\Delta_1$, so only asphericity of all products $I\mathrm{x}a$ remains to check. The reference to the cancelled corollary is revised to the end of section~36. Page~64 literally says ``for a in J'' where the construction indexes by $a\in A$, and retains a stray typed $u$ beneath a final-object caret; these visible irregularities remain in the readers. Section~39 returns to total asphericity and refines the terminology accordingly.

\subsection*{39. Strict, local and wide-sense test categories (pp. 66--69)}
The return to total asphericity first reuses the homotopy lemma: if every representable admits a contraction to a chosen section through a fixed homotopy interval, then every representable is aspheric over the final object and $\widehat A$ is totally aspheric (p.~67). The proof does not use separation of the endpoints. The following assertion that separation follows from the other assumptions is too broad if those assumptions are taken alone: the terminal category with $I=e$ supplies a counterexample. The boxed and diagonally cancelled paragraph had discussed this exceptional case; its removal is preserved, together with the still-live preceding assertion. This is one of the places where the revision marks matter to the logical reading.

The terminology is now reorganized around the Lawvere element. A local test category has $L_A$ aspheric over the final presheaf; a test category is additionally aspheric; a strict test category additionally has a totally aspheric presheaf topos. Their presheaf categories are correspondingly local elementary modelizers, elementary modelizers and elementary strict modelizers. A general strict modelizer is characterized by its weak-equivalence class and a localization equivalent to $(\mathrm{Hot})$ that preserves finite products; further requirements concerning sums and infinite products are still being considered (pp.~67--68). The weaker adjunction criterion $i_Aj_A(C)\to C$ remains the definition of a test category ``in the wide sense''. The source's dismissal of the likely usefulness of that condition is an uncancelled parenthetical, while an adjacent faint marginal note is more favourable; its exact wording and final Delta superscript remain uncertain. The page's note call is 45. The later shorthand ``T1 + T3'' must be read with care: the earlier T3 on p.~39 asked only for an aspheric object with disjoint endpoints, and its margin used T2 to pass to asphericity over the final object. The operative interval condition in the present definition is explicitly relative. The terminology for general modelizing topoi is proposed with an express reservation that its connection to actual models has not yet been established (p.~68).

The distinction between test and strict test categories is then made concrete. A first proposed small category is generated by interval powers together with $a_o\mathrm{x}I^n$; Grothendieck expects the missing product $a_o\mathrm{x}a_o$ in the presheaf category to fail asphericity. A more decisive construction takes any test category $A$ and a homotopy interval $I$: the slice $A_{/I}$ is again local test and aspheric, hence a test category, but is never strict because the two endpoint subobjects of its final object are nonempty and disjoint (p.~69). Total asphericity would force their intersection to be aspheric, so it cannot hold. The first construction remains an expectation, whereas the slice construction supplies the stated universal counterexample. The scan's opening insertion on p.~69 reads ``as'', giving ``and as maps \ldots\ all those'', rather than the witness-based ``those'' of the initial transcription. Section~40 next supplies the product-stability lemma that the preceding arguments have been using.

\subsection*{40. Products and relative weak equivalences (pp. 69--71)}
The product argument supplies a fact already used in section~37: if $A$ is aspheric, the projection $C\mathrm{x}A\to C$ is aspheric for every $C$. The slice criterion reduces first to $C$ with a final element and then to products $C\mathrm{x}A_{/a}$, which themselves have final elements (p.~69). Grothendieck strengthens this to stability of weak equivalences under cartesian products. It is enough to treat $f\mathrm{x}\mathrm{id}_C$ with $f:A'\to A$ a weak equivalence. Viewing this as a morphism over $C$ motivates weak equivalence relative to a base topos $Y$: the property must survive each localization by an object of $Y$, and testing on generators suffices. For categories over $C$, the test is the induced functor over every $c\in C$. Thus the product statement reduces to $C$ with a final element, where the projection square and two-out-of-three prove it (p.~70). The displayed localization arrow literally has target $U$ in $Y_{/U}\to U$; its described role requires target $Y$, so the reader preserves the visible letter and flags the mismatch here. The final square's upper handwritten label has lowercase $c$, another source detail retained as drawn.

Universal weak equivalences strengthen this relative notion by requiring stability under every base change in $(\mathrm{Cat})$, not only localizations (p.~71). Their test on base-change categories with a final element suggests an analogy with trivial fibrations. Grothendieck proposes that fibering or cofibering functors with aspheric fibers, and more generally such functors whose transition functors are weak equivalences, might supply the categorical counterparts of Quillen's fibrations. The inserted qualification ``in the purely categorical sense'' limits the proposal. Cofibrations remain unidentified, but monomorphisms alone appear insufficient, prompting the revised references to section~30 and (9). These are prospective comparisons rather than an established model structure on $(\mathrm{Cat})$. The margin labels the terminology ``awkward'' and carries call (46). The section then deliberately returns to test categories, leaving section~41 to clarify exactly which assumptions the preceding proofs have used.

\subsection*{41. The guiding assumption and weakened saturation (pp. 71--72)}
The conjectural triviality of the category of autoequivalences of $(\mathrm{Hot})$ had encouraged defining a modelizer through its weak-equivalence class, with the functor to homotopy types following essentially uniquely. Grothendieck now emphasizes, by changing the final period to an exclamation mark, that none of the proved statements has used that guiding assumption (p.~71). The operative definition of modelizer has also changed: the initial requirement of strong saturation, meaning that $W$ consists of precisely the arrows inverted by localization, has been replaced by the three weaker conditions written in section~37. The authorial transposition under ``wrote explicitly down'' supplies ``wrote down explicitly'' in the applied reading; the diplomatic text keeps the original order and the pen curve. The source spelling ``drapping'' is left visible where the earlier transcription had conjectured an $o$ correction.

Strong saturation can nevertheless be transferred between $(\mathrm{Cat})$ and an elementary modelizer through the established comparison functors, so one example suffices for all (p.~72). Grothendieck invokes the known closed model structure on simplicial sets and his cited proposition~1 in Quillen's I~5.5 for that example. His remaining question concerns the proof of that model structure: can the present methods replace the specialized simplicial argument by one applicable to general test categories, or at least strict ones? He explicitly postpones this question and returns to comparing test categories and completing the inventory of standard examples. The handwritten call following ``strongly saturated'' is (47), while the separate margin has an asterisk and a reference to underlined VIII whose short abbreviation remains uncertain. This transition motivates section~42's matching revision of the definition of a model-preserving functor.

\subsection*{42. Model-preserving functors and full subcategories (pp. 72--74)}
The weakened saturation convention of section~41 requires a corresponding change to the definition of model-preserving functor: in addition to inducing an equivalence after localization, $f:M\to M'$ must satisfy the equality $W=f^{-1}(W')$, so that it reflects as well as preserves the specified weak equivalences (p.~72). The revised text first returns to the unresolved modelizing status of ordered sets and preorders. It then distinguishes the equivalence of localized categories from a canonical comparison with the chosen category of homotopy types. The earlier rigidity assumption would supply such a comparison, but has not been established. A subsidiary question asks whether an endofunctor of $(\mathrm{Hot})$ that is the identity on objects must be naturally isomorphic to the identity (p.~73). These questions remain conjectural; they do not become consequences merely from the existence of the concrete comparison functors already constructed.

The constructive result uses a full subcategory containing the models produced by a test category. In the preliminary abstract argument an endofunctor $h$ of a modelizer factors through $M'\subset M$ and induces an autoequivalence on the localization; the text then describes the induced inclusion and factor as quasi-inverse (pp.~73--74). That exact formulation is stronger than its stated hypothesis: their composite is the localized $h$, and an autoequivalence need not be isomorphic to the identity. For example, taking $M'=M$ and inclusion equal to the identity still leaves that distinction. In the intended application $h=i_Ai^*$, the key comparison theorem does provide the required natural comparison to the identity. The resulting proposition says that a full subcategory of $(\mathrm{Cat})$ containing all categories of elements $A_{/F}$ for a test category $A$ inherits a modelizer structure and has model-preserving inclusion (p.~74). Categories locally isomorphic to $A$ supply a useful reformulation of the required objects. This reopens the ordered-set problem: a suitable test category defined by an ordered set would put all its categories of elements inside $(\mathrm{Ord})$. Such a category could not be strict, by the earlier obstruction for topoi attached to topological spaces in section~35. The last remarks contrast localization stability for test categories with the stricter total-asphericity condition, and lead directly to the standard examples of section~43. The source retains its literal ``T1 + T3'' shorthand. Its interlinear ``inclusion'', ``modelizer'', ``the test category'', ``algebraic'' and ``a suitable'' are typed revisions; the cancelled locally-ordered qualification is a genuine pen deletion.

\subsection*{43. Ordered, unordered and face-only simplices (pp. 75--78)}
The return to standard examples applies section~38's test-functor criterion directly. For ordered simplices, the inclusion $i:\Delta\to(\mathrm{Cat})$ is fully faithful and carries the interval to the ordinary two-object ordered category. Thus $i^*(\Delta_1)$ is the representable interval, and the previously proved asphericity of its products with simplices gives the test condition. Contractions to final vertices then give total asphericity and hence strictness (p.~75). The cancelled preliminary proof is replaced by this more precise argument. For unordered simplices, a small skeleton of nonempty finite sets has finite products and a two-point interval. Sending a finite set to its ordered set of nonempty subsets supplies the test functor. On the two-point set this produces the span of the two singleton subsets into the full set; the endpoint-preserving map to $\Delta_1$ sends $\{0\}$ to 0 and both other subsets to 1 (pp.~75--76). The comparison square with ordered sets identifies the resulting functor on ordered simplices with the flag, or barycentric-subdivision, construction. The single untilded $\Delta$ in the unordered-category conclusion on p.~75 is a visible source omission; the later functor-domain notation has an additional uncertain wavy stroke, rather than a securely readable presheaf hat.

Removing degeneracies gives a substantially different answer. Grothendieck first emphasizes that face operators already determine homology, cohomology and the fundamental groupoid, and hence the cohomological test for a weak equivalence of the associated complexes. Nevertheless, in the category $\Delta^f$ of strictly increasing simplex maps, the product of the zero-simplex with an $n$-simplex has a discrete category of elements with $n+1$ objects. It is therefore not aspheric when $n>0$ (p.~76). Removing the zero-simplex does not supply a representable interval with sections, since arbitrarily large simplices cannot map into a fixed smaller one. The more flexible attempt using the presheaf $i^*(\Delta_1)$ also fails the required relative interval test. The canonical Lawvere element gives the same two-point obstruction over the zero-simplex (p.~77). Thus $\Delta^f$ is not a local test category, while the possibility that it is a test category in the wide sense remains open at this point. The added name ``weak test category'' anticipates the next section's systematic terminology.

The remaining weak condition asks for absolute asphericity of $i_Aj_A(C)$ for every category $C$ with a final object. Its first tests are asphericity of $A$ itself and of the Lawvere element, expressly distinguished from asphericity over the final presheaf (p.~77). A positive answer would let section~42's proposition construct further full-subcategory modelizers. To pursue the ordered-set case, Grothendieck considers the poset $J$ of nonempty finite subsets of an infinite vertex set. It has the desired local simplex structure and, for countable vertices, is a category of elements of an infinite face simplex. Its filtering property suggests absolute asphericity. The local test condition still fails because a minimal vertex has a one-point slice, which cannot be a test category (p.~78). The closing calculation identifies the category of elements of $L_J$ with pairs $K\subset T$, with $T$ finite and nonempty, ordered by $T'\subset T$ and $K'=K\cap T'$. It begins an explicit test of the surviving weak condition rather than asserting its conclusion. Section~44 opens the following day by announcing that the face-only simplex category does satisfy that weaker notion. The readers retain the source's shifted apostrophes, the cancelled product/intersection argument on p.~78, and the distinction between the two inline typed round-letter indices and the displayed handwritten infinity.

\subsection*{44. Weak, local, strict test categories and test functors (pp. 79--88)}
The discovery that the face-only simplex category is a weak test category rehabilitates the weakest member of the three notions and prompts a systematic comparison. Part~a recalls the adjunction $i_A:\widehat A\rightleftarrows(\mathrm{Cat}):j_A$, with $i_A(F)=A_{/F}$, and compares eight formulations of the weak test condition: equivalence of localized categories, control of the adjunction maps, asphericity of $A_{/j_A(C)}$ for $C$ with a final object, and preservation of weak equivalences by $i_Aj_A$ together with asphericity of $A$ (pp.~79--81). The local asphericity argument identifies the stronger counit conditions; the remaining implication from the initial equivalence condition is revisited in the later marginal note dated 27.3, which observes that every autoequivalence preserves a final object. Thus the extra final-object hypothesis is unnecessary in that formulation. The notes numbered 48 and 49 have broad pencil hatching whose cancellation status remains doubtful. A literal composition on p.~80 is reversed: condition~(v) prints $j_Ai_A(C)$ where both its domain and the displayed category of elements require $i_Aj_A(C)$. This source slip is preserved. Part~b identifies local test categories through weak test slices, a separated interval aspheric over the final presheaf, the corresponding property of the Lawvere element, and asphericity of the comparison $i_Aj_A(C)\to A\times C$ (pp.~81--83). Its proof uses the interval comparison and homotopy argument of section~37. At the reduction on p.~82 the text says that $A$ has a final element, although replacing $A$ by $A_{/F}$ for aspheric $F$ supplies only asphericity in general; the argument needs that weaker condition. The same page's final reference says ``in same \S.31'', despite the preceding reference to section~37. These formulations remain in the readers. A local test category that is also aspheric is a test category. Grothendieck calls the corresponding presheaf topoi locally modelizing and modelizing, and asks whether the weak condition can also be made intrinsic to the topos rather than to its chosen generating category. Part~c characterizes strictness by preservation of binary products after localization, equivalently by total asphericity together with the local test property (pp.~83--85). The concise conditions T~2 and T~3 require the added hypothesis $A\ne\varnothing$. Their equivalence with the product formulation uses the asphericity test for $A_{/F\times G}\to A_{/F}\times A_{/G}$. The source's reference call is 36, and the interlinear words ``in A'' with their added hat do not cancel the following word ``of'' on p.~84. Part~d introduces weak test functors $i:A\to(\mathrm{Cat})$ through the model-preserving property of $i^*$, including reflection of weak equivalences; equivalently the composite $i_Ai^*$ is model preserving (pp.~85--86). An induced autoequivalence alone supplies no chosen comparison with the identity. With final objects in the $i(a)$, the concrete evaluation functor $A_{/i^*(C)}\to C$ gives that comparison by the earlier homotopy argument. The revised proposition on p.~86 states only $(ii)\Leftrightarrow(iii)\Rightarrow(i)$: relative asphericity of $i^*(\Delta_1)$ is equivalent to an endpoint-compatible interval map $i_!(I)\to\Delta_1$, with asphericity of $A$ explicitly added. It does not assert the converse from weak test functor to these stronger conditions. The latter define a test functor on p.~87. Whether every weak test functor on a test category is of this stronger sort remains open here; strictness of $A$ removes the absolute-versus-relative asphericity obstruction. The closing remarks then correct a further logical shortcut: recovering reflection of weak equivalences from an induced equivalence uses strong saturation of $W$, and the necessary Quillen-model argument is not yet checked in this presentation (p.~88). The left adjoint $i_!$ need not itself be model preserving, even for a strict test category. Grothendieck also leaves open whether model preservation of $i_A$, or of a more general $i^*$, already forces the weak test property; the canonical $j_A$ with asphericity is the known case. The deleted fourth remark returns to the final-object observation already incorporated in the earlier margin. The arXiv witness omits the live ``even'' in Remark~2 and rewrites the ending of Remark~3, losing ``assuming'', ``is'', and its comma; these omissions obscure the successive qualifications preserved here. The next main section uses this review to examine the interplay of weak equivalences, aspheric objects, homotopy intervals and test functors; the intervening archive leaves contain later notes to this chapter.

\subsection*{Later notes to Chapter II (cote 134-3, sites 89--93; leaves II i--v)}
These five leaves follow the main chapter in archive order and annotate it retrospectively; their lower foliation 89--93 must not be confused with the next chapter's main pages 89--93 in cote 134-4. Notes~1--3 revisit the Hauptvermutung, the semisimplicial construction of the Leray spectral sequence, and the attribution of the Dold--Kan theorem. Note~3 expressly calls ``Dold-Puppe'' an inaccurate name, contrary to the witness's ``accurate''. Note~4 retracts the earlier obstruction to expressing stack cohomology through the primitive source-and-target structure, introduces hemispherical or globular complexes, and identifies the category of standard hemispheres with the one used in David W.~Jones's thesis. Its face operations have indices $i\geq1$, whereas its degeneracies begin at $i\geq0$; the arXiv witness changes the latter bound to 1. The handwritten category symbol has an inner curved arc, retained in the readers. The note compares hemispheres, simplices and cubes as test categories and contractors, then asks for explicit test functors and suitable existence and uniqueness results, pointing to section~90 (sites 89--90). Notes~5--6 refer to the Breen appendix. Note~7 recalls the nerve construction in a Bourbaki talk but leaves the bibliographical slot blank, marked by a marginal X; that unfinished reference is not supplied editorially. Notes~8--9 retract the proposed topology characterization and the early cofibration description, directing the reader to sections~24 and~30 and to the later discussion of Thomason in section~87. Notes~10--12 record the expanding scope of the project, the preference for hemispheres, and the connection with Benabou's work (site~90). Note~13 makes the larger methodological change explicit: Quillen's axioms recede as asphericity and contractibility structures become the more immediate tools for constructing modelizers and their canonical homotopy-type functors. Notes~15 and~19 reject the unrestricted use of the canonical topology, especially for $(\mathrm{Cat})$. Note~16 bases the suitability of mixed multiple complexes on the test-category criterion and stability under products, referring to sections~31 and~74. Note~17 corrects the earlier compactness description of cells: a finiteness condition on the ordered set is needed. Notes~18 and~20 leave the relevant modelizer assertion unproved and set finite-presentation restrictions aside (site~91). Notes~21--22-prime describe the displacement of cellular structures by integrators and cointegrators, and of homotopy-equivalent topologies by test notions; the proper-map correction in note~22 must be read with the condition it qualifies in the earlier body. Notes~23--28 identify the slice-topos equivalence, the evolving test terminology, canonical modelizers, strictly generating families, and the eventual abandonment of topological realization as the main route. The fully handwritten note~28-prime points to the weakened condition of section~41. Note~29 still treats the modelizing status of ordered sets and preorders as unproved. Notes~30--32 connect product test categories and their natural comparison functors with sections~74 and~82, prefer hemispheres again, and question the efficiency of the semisimplicial route to the Cech--Cartier--Verdier functor (site~92). The last leaf continues that question toward a projected return to topoi, revises a section reference from 38 to 39, and cancels an entire example about failure of total asphericity in slices. Note~35 explicitly retracts the earlier confusion between ordinary limits and homotopy limits. Its claim that only finite products and arbitrary sums exist is itself too restrictive on products: arbitrary products exist in the homotopy category, represented by products of Kan complexes, since maps and homotopies into a product are determined componentwise. The source wording is preserved. Note~36 visibly has a corrected numeral XIII, with an uncertain older layer. The dated note~37 strengthens a modelizing condition by also requiring $j_A=i_A^*$ to be modelizing and refers to section~65, D), Remark~3. Note~38 finally sends the reader back to section~44 for the revised weak, ordinary and strict test notions. Throughout these leaves the checked text distinguishes typed interlinear amendments from handwriting; in particular the live final i and typed ve of ``primitive'' must survive, while the apparent commas attached to several insertion hooks are not textual punctuation.

\subsection*{45. Comparing modelizers by test functors (pp. 89--92)}
The review of section~44 leads back to the organizing problem of section~26: how weak equivalences, aspheric objects, homotopy intervals and test functors determine one another in a general modelizer (p.~89). Test categories are compared with spectra of discrete valuation rings, which can be sent into other objects to test their structure. A proposed test functor $i:A\to M$ should induce a modelizing functor $i^*:M\to\widehat A$, whereas the canonical $i_A:\widehat A\to(\mathrm{Cat})$ sends an elementary modelizer toward the common categorical model (p.~90). Their exactness properties run in opposite directions: $i^*$ preserves inverse limits, while $i_A$ preserves direct limits. The left adjoint $i_!$ is generally not expected to be model preserving; the canonical pair $(i_A,j_A)$ is exceptional in having both adjoints model preserving. Composing gives $i_Ai^*:M\to(\mathrm{Cat})$, but this composite generally loses both exactness properties. For a strict modelizer its preservation of finite products survives up to weak equivalence. The discussion also points to preservation of sums, with a faint pencil margin adding fibered products; that addition is absent from the witness. In the reverse direction, finding $(\mathrm{Cat})\to M$ is equivalent, through the comparison functors already available, to finding a modelizing functor $\widehat A\to M$. Neither existence nor useful exactness is asserted for arbitrary $M$. Topological spaces provide a familiar special case: a suitable realization $r:A\to(\mathrm{Spaces})$ extends to $r_!$, and the adjunction $(r_!,r^*)$ can have both members modelizing (p.~91). The resulting categorical realization is explicitly $r_!i^*:(\mathrm{Cat})\to(\mathrm{Spaces})$ in formula~(6), not the witness's $r_!r^*$, whose right-hand factor would have the wrong source here. This analogy does not make spaces an equally satisfactory basic modelizer: their structure is less algebraic, the realization choices are less canonical, and the existence of such adjoint comparisons for broad classes of test categories remains unsettled (p.~92). The conclusion therefore singles out the two constructions $i^*$ and $i_A$ as the main ones and prepares the definitions in section~46. The witness omits the intervening words ``test functor (still to be defined !) of a test category into an arbitrary'' from the first construction's description on p.~92; retaining them is necessary to see which functor remains to be defined. The checked source also retains the author's paragraph instruction before the final reflection and the joining curve in ``def ined'', without converting either into manuscript prose.

\subsection*{46. Weak test functors into an arbitrary modelizer (pp. 93--95)}
A modelizer here consists of a category $M$ and a class $W_M$ satisfying the mild saturation conditions of p.~59, with localization equivalent to $(\mathrm{Hot})$; no particular equivalence is included in the structure. The ``inspiring assumption'' supplies motivation but is expressly withheld from the technical argument (p.~93). For a weak test category $A$, the proposed weak test condition on $i:A\to M$ is that $i^*:M\to\widehat A$ be modelizing: it must reflect as well as preserve weak equivalences and induce an equivalence of localizations. The source uses $W_A$ in this formula, without a hat on its subscript. Aspheric objects are defined by their maps to a final object $e_M$. Their images are final in $H_M$ if $e_M$ remains final there; the converse additionally uses strong saturation. The dated marginal note 29.3 observes that existence of a weak test functor verifies the final-object condition. In the special categorical case, the striking extra feature is an equivalence between preservation of weak equivalences, preservation of aspheric objects and the full weak test condition (p.~94). The explanation explicitly restores a hypothesis omitted from the displayed observation: every category $i(a)$ must have a final object. With this hypothesis, section~37's key result on p.~61 proves the nontrivial implication, and it even suffices to test categories $C$ having final objects. The inserted reference is $(a\,iv)\Rightarrow(a\,ii)$, with typed clauses and a handwritten arrow. The resulting question asks when a functor between general modelizers that preserves weak equivalences, or merely aspheric objects, must already be model preserving. Possible assumptions include compatibility with inverse limits or existence of a left adjoint, but no general theorem is asserted (pp.~94--95). The first proposed test cases are $i^*:M\to\widehat A$ with $M=(\mathrm{Cat})$ or another elementary modelizer $\widehat B$. They lead directly to the deferred task of dropping the final-object hypothesis on $i(a)$. The following section begins on 30.3 by announcing a satisfactory answer under strictness of $A$; the explicit words ``test category'' in that announcement are omitted by the witness. Several dark marginal or punctuation-like marks remain unresolved in the checked source, including the raised mark after $M$ in the general question; they do not supply additional hypotheses.

\subsection*{47. The categorical comparison through functor categories (pp. 95--98)}
The proposed extension starts with a distinction that the notation must preserve. The presheaf $i^*(C)$ has sets $\mathrm{Hom}(i(a),C)$ as values, whereas the category $A_{//C}$ is fibered over $A$ with the full functor categories $\underline{\mathrm{Hom}}(i(a),C)$ as fibers (p.~96). The set of objects, regarded as a discrete category, gives a canonical comparison $A_{/C}\to A_{//C}$, usually far from an equivalence on individual fibers. Constant functors supply the other comparison, giving the zigzag $A_{/C}\to A_{//C}\leftarrow A\times C\to C$. This is the first explicit use of the bicategory structure of $(\mathrm{Cat})$ in the argument: natural transformations now provide morphisms inside the fibers. The statement that $C$ is identified with the full subcategory of constant functors needs qualification before assumptions on $i(a)$ are imposed. It is valid for nonempty connected $i(a)$; with several connected components, transformations between constant functors can vary by component, and an empty domain collapses the objects. The later contractibility condition supplies the needed connectedness and nonemptiness. If all three maps in the zigzag are weak equivalences, they give a canonical isomorphism between $i_Ai^*(C)$ and $C$ after localization. Thus $i_Ai^*$ preserves weak equivalences and induces the identity up to isomorphism. When $A$ is weak test, the known equivalence induced by $i_A$ then transfers the equivalence property to $i^*$ (p.~97). The source calls the latter a ``test-functor'' at this point, although the earlier definition applies that name to $i$ and calls $i^*$ modelizing.

The three arrows lead to different hypotheses. The projection $A\times C\to C$ is a weak equivalence for all $C$ exactly when $A$ is aspheric, by section~40. For the middle arrow, Grothendieck admits a general fiberwise comparison proposition: a cartesian functor between fibered categories, or a cocartesian functor between cofibered categories, is a weak equivalence if it is so on every fiber. Its second alternative literally prints ``i is cartesian'', although the proposition quantifies the functor $u$; that variable slip remains in the source. This reduces the middle arrow to the maps $C\to\underline{\mathrm{Hom}}(i(a),C)$. Their being weak equivalences for every $C$ is announced as a characterization of contractibility of $i(a)$, a stronger notion than mere asphericity. The marginal correction dated 9.4 postpones the promised proof to p.~121 and its converse to p.~143. The first arrow is then expected to yield to the asphericity criterion, producing a generalized key result with twice as many equivalent formulations (p.~98). This section states the program and its reductions; it does not yet prove that final equivalence theorem. The closing terminology correction explains why earlier uses of ``contractible'' for ``aspheric'' must be read cautiously outside the CW-space setting. Categories with final objects remain the evident contractible examples. The next section develops homotopies and contractibility in a general category with weak equivalences. Its opening retains the explicit refusal to assume a strict initial object until needed, contrary to the witness's omission of ``not'', and adds binary products as a precaution.
